\batchmode


\documentclass[12pt,spanish]{article}
\RequirePackage{ifthen}



\usepackage[utf8]{inputenc}
\usepackage[spanish]{babel}

\usepackage{framed}


\usepackage{fancyhdr}

\usepackage{anysize} 

\usepackage{dsfont}	

\usepackage[dvips]{graphicx}

\DeclareGraphicsExtensions{.pdf,.png,.jpg}  

\usepackage{color}


\usepackage{multicol}\definecolor{gray97}{gray}{.97}
\definecolor{gray75}{gray}{.75}
\definecolor{gray45}{gray}{.45}

\usepackage{listings}
\usepackage[hidelinks]{hyperref}
\usepackage{setspace}


\lstdefinelanguage[]{omi} { 
   %
frame=Ltb,
     framerule=0pt,
     aboveskip=0.2cm,
     framextopmargin=0.3pt,
     framexbottommargin=0.2pt,
     framexleftmargin=0.5cm,
     framesep=0pt,
     rulesep=.1pt,
     tabsize=2,
     backgroundcolor=\color{gray97},
     rulesepcolor=\color{black},
     %
stringstyle=\ttfamily ,
     showstringspaces = false,
     basicstyle=\scriptsize\ttfamily ,
     commentstyle=\color{gray75},
     comment=[l]{\#},%
morecomment=[l]{//},%
morecomment=[s][\color{gray75}]{/*}{*/},%
morecomment=[s][\color{green}]{`}{`},%
keywordstyle=\bfseries ,
     stringstyle=\itshape\ttfamily\color{red},
     string=[b]',%
morestring=[b]",%
escapeinside=||,
   	%
numbers=left,
     numbersep=1pt,
     numberstyle=\tiny ,
     numberfirstline = false,
     breaklines=true,
     emph={%
$PS1, $PS2, ~, switch, case, default, goto, as, in, function, for, while, false, true, size, class, extends, parent, \<\<, private, static, return, &, include, global,~,new, null, if, else, this, $, prompt$, elif, break, continue, load %
},
     emphstyle={\bfseries\itshape\sffamily },%
}%
\lstdefinelanguage[]{bash} { 
   %
frame=Ltb,
     framerule=0pt,
     aboveskip=0.2cm,
     framextopmargin=0.3pt,
     framexbottommargin=0.2pt,
     framexleftmargin=0.5cm,
     framesep=0pt,
     rulesep=.1pt,
     tabsize=2,
     backgroundcolor=\color{gray97},
     rulesepcolor=\color{black},
     %
stringstyle=\ttfamily ,
     showstringspaces = false,
     basicstyle=\scriptsize\ttfamily ,
     commentstyle=\color{gray75},
     comment=[l]{\#},%
morecomment=[s][\color{gray75}]{/*}{*/},%
keywordstyle=\bfseries ,
     stringstyle=\itshape\ttfamily\color{red},
   	%
numbers=left,
     numbersep=1pt,
     numberstyle=\tiny ,
     numberfirstline = false,
     breaklines=true,
     emph={%
$PS1, $PS2, std, c++11, g++ %
},
     emphstyle={\bfseries\itshape\sffamily },%
}%
\lstdefinelanguage[]{cpp} { 
   %
frame=Ltb,
     framerule=0pt,
     aboveskip=0.2cm,
     framextopmargin=0.3pt,
     framexbottommargin=0.2pt,
     framexleftmargin=0.5cm,
     framesep=0pt,
     rulesep=.1pt,
     tabsize=2,
     backgroundcolor=\color{gray97},
     rulesepcolor=\color{black},
     %
stringstyle=\ttfamily ,
     showstringspaces = false,
     basicstyle=\scriptsize\ttfamily ,
     commentstyle=\color{gray75},
     comment=[l]{//},%
morecomment=[s][\color{gray75}]{/*}{*/},%
morecomment=[s][\color{green}]{<omi}{h>},%
morecomment=[s][\color{red}]{"}{"},%
keywordstyle=\bfseries ,
     stringstyle=\itshape\ttfamily\color{red},
   	%
numbers=left,
     numbersep=1pt,
     numberstyle=\tiny ,
     numberfirstline = false,
     breaklines=true,
     emph={%
class,include, define, public, void, this, if, else, throw, extern, using
     },
     emphstyle={\bfseries\itshape\sffamily },%
}

\usepackage{makeidx}
\makeindex 

\parskip 0.2in									
%
\newenvironment{changemargin}[2]{%
\begin{list}{}{%
\setlength{\topsep}{300pt}%
\setlength{\leftmargin}{#1}%
\setlength{\rightmargin}{#2}%
\setlength{\listparindent}{\parindent}%
\setlength{\itemindent}{\parindent}%
\setlength{\parsep}{\parskip}%
}%
\item[]}{\end{list}} \usepackage{float}
\restylefloat{table}
\usepackage{placeins}
\usepackage{framed}
\usepackage{epsf}

\usepackage{titlesec}


\setcounter{secnumdepth}{4}
\setcounter{tocdepth}{5}


\titleformat{\paragraph}
{\normalfont\normalsize\bfseries }{\theparagraph }{1em}{}
\titlespacing*{\paragraph}
{0pt}{3.25ex plus 1ex minus .2ex}{1.5ex plus .2ex}


\usepackage{amsthm}
\usepackage{amsmath}
\usepackage{array}
\usepackage{eurosym} 
\usepackage{fancybox}
\usepackage{float}
\usepackage[T1]{fontenc}


\author{Manual de usuario \\\\\ Fco. Javier Bohórquez Ogalla}						%
\date{}														%


\title{ 
\begin{center}
\includegraphics[scale=0.5]{logo-doc.png}
\end{center} 
}


\DeclareUnicodeCharacter{E6}{\ae}
\usepackage{dingbat}




\pagecolor[gray]{.7}

\usepackage[latin1]{inputenc}



\makeatletter

\makeatletter
\count@=\the\catcode`\_ \catcode`\_=8 
\newenvironment{tex2html_wrap}{}{}%
\catcode`\<=12\catcode`\_=\count@
\newcommand{\providedcommand}[1]{\expandafter\providecommand\csname #1\endcsname}%
\newcommand{\renewedcommand}[1]{\expandafter\providecommand\csname #1\endcsname{}%
  \expandafter\renewcommand\csname #1\endcsname}%
\newcommand{\newedenvironment}[1]{\newenvironment{#1}{}{}\renewenvironment{#1}}%
\let\newedcommand\renewedcommand
\let\renewedenvironment\newedenvironment
\makeatother
\let\mathon=$
\let\mathoff=$
\ifx\AtBeginDocument\undefined \newcommand{\AtBeginDocument}[1]{}\fi
\newbox\sizebox
\setlength{\hoffset}{0pt}\setlength{\voffset}{0pt}
\addtolength{\textheight}{\footskip}\setlength{\footskip}{0pt}
\addtolength{\textheight}{\topmargin}\setlength{\topmargin}{0pt}
\addtolength{\textheight}{\headheight}\setlength{\headheight}{0pt}
\addtolength{\textheight}{\headsep}\setlength{\headsep}{0pt}
\setlength{\textwidth}{349pt}
\newwrite\lthtmlwrite
\makeatletter
\let\realnormalsize=\normalsize
\global\topskip=2sp
\def\preveqno{}\let\real@float=\@float \let\realend@float=\end@float
\def\@float{\let\@savefreelist\@freelist\real@float}
\def\liih@math{\ifmmode$\else\bad@math\fi}
\def\end@float{\realend@float\global\let\@freelist\@savefreelist}
\let\real@dbflt=\@dbflt \let\end@dblfloat=\end@float
\let\@largefloatcheck=\relax
\let\if@boxedmulticols=\iftrue
\def\@dbflt{\let\@savefreelist\@freelist\real@dbflt}
\def\adjustnormalsize{\def\normalsize{\mathsurround=0pt \realnormalsize
 \parindent=0pt\abovedisplayskip=0pt\belowdisplayskip=0pt}%
 \def\phantompar{\csname par\endcsname}\normalsize}%
\def\lthtmltypeout#1{{\let\protect\string \immediate\write\lthtmlwrite{#1}}}%
\newcommand\lthtmlhboxmathA{\adjustnormalsize\setbox\sizebox=\hbox\bgroup\kern.05em }%
\newcommand\lthtmlhboxmathB{\adjustnormalsize\setbox\sizebox=\hbox to\hsize\bgroup\hfill }%
\newcommand\lthtmlvboxmathA{\adjustnormalsize\setbox\sizebox=\vbox\bgroup %
 \let\ifinner=\iffalse \let\)\liih@math }%
\newcommand\lthtmlboxmathZ{\@next\next\@currlist{}{\def\next{\voidb@x}}%
 \expandafter\box\next\egroup}%
\newcommand\lthtmlmathtype[1]{\gdef\lthtmlmathenv{#1}}%
\newcommand\lthtmllogmath{\dimen0\ht\sizebox \advance\dimen0\dp\sizebox
  \ifdim\dimen0>.95\vsize
   \lthtmltypeout{%
*** image for \lthtmlmathenv\space is too tall at \the\dimen0, reducing to .95 vsize ***}%
   \ht\sizebox.95\vsize \dp\sizebox\z@ \fi
  \lthtmltypeout{l2hSize %
:\lthtmlmathenv:\the\ht\sizebox::\the\dp\sizebox::\the\wd\sizebox.\preveqno}}%
\newcommand\lthtmlfigureA[1]{\let\@savefreelist\@freelist
       \lthtmlmathtype{#1}\lthtmlvboxmathA}%
\newcommand\lthtmlpictureA{\bgroup\catcode`\_=8 \lthtmlpictureB}%
\newcommand\lthtmlpictureB[1]{\lthtmlmathtype{#1}\egroup
       \let\@savefreelist\@freelist \lthtmlhboxmathB}%
\newcommand\lthtmlpictureZ[1]{\hfill\lthtmlfigureZ}%
\newcommand\lthtmlfigureZ{\lthtmlboxmathZ\lthtmllogmath\copy\sizebox
       \global\let\@freelist\@savefreelist}%
\newcommand\lthtmldisplayA{\bgroup\catcode`\_=8 \lthtmldisplayAi}%
\newcommand\lthtmldisplayAi[1]{\lthtmlmathtype{#1}\egroup\lthtmlvboxmathA}%
\newcommand\lthtmldisplayB[1]{\edef\preveqno{(\theequation)}%
  \lthtmldisplayA{#1}\let\@eqnnum\relax}%
\newcommand\lthtmldisplayZ{\lthtmlboxmathZ\lthtmllogmath\lthtmlsetmath}%
\newcommand\lthtmlinlinemathA{\bgroup\catcode`\_=8 \lthtmlinlinemathB}
\newcommand\lthtmlinlinemathB[1]{\lthtmlmathtype{#1}\egroup\lthtmlhboxmathA
  \vrule height1.5ex width0pt }%
\newcommand\lthtmlinlineA{\bgroup\catcode`\_=8 \lthtmlinlineB}%
\newcommand\lthtmlinlineB[1]{\lthtmlmathtype{#1}\egroup\lthtmlhboxmathA}%
\newcommand\lthtmlinlineZ{\egroup\expandafter\ifdim\dp\sizebox>0pt %
  \expandafter\centerinlinemath\fi\lthtmllogmath\lthtmlsetinline}
\newcommand\lthtmlinlinemathZ{\egroup\expandafter\ifdim\dp\sizebox>0pt %
  \expandafter\centerinlinemath\fi\lthtmllogmath\lthtmlsetmath}
\newcommand\lthtmlindisplaymathZ{\egroup %
  \centerinlinemath\lthtmllogmath\lthtmlsetmath}
\def\lthtmlsetinline{\hbox{\vrule width.1em \vtop{\vbox{%
  \kern.1em\copy\sizebox}\ifdim\dp\sizebox>0pt\kern.1em\else\kern.3pt\fi
  \ifdim\hsize>\wd\sizebox \hrule depth1pt\fi}}}
\def\lthtmlsetmath{\hbox{\vrule width.1em\kern-.05em\vtop{\vbox{%
  \kern.1em\kern0.8 pt\hbox{\hglue.17em\copy\sizebox\hglue0.8 pt}}\kern.3pt%
  \ifdim\dp\sizebox>0pt\kern.1em\fi \kern0.8 pt%
  \ifdim\hsize>\wd\sizebox \hrule depth1pt\fi}}}
\def\centerinlinemath{%
  \dimen1=\ifdim\ht\sizebox<\dp\sizebox \dp\sizebox\else\ht\sizebox\fi
  \advance\dimen1by.5pt \vrule width0pt height\dimen1 depth\dimen1 
 \dp\sizebox=\dimen1\ht\sizebox=\dimen1\relax}

\def\lthtmlcheckvsize{\ifdim\ht\sizebox<\vsize 
  \ifdim\wd\sizebox<\hsize\expandafter\hfill\fi \expandafter\vfill
  \else\expandafter\vss\fi}%
\providecommand{\selectlanguage}[1]{}%
\makeatletter \tracingstats = 1 
\providecommand{\Chi}{\textrm{X}}
\providecommand{\Epsilon}{\textrm{E}}
\providecommand{\Rho}{\textrm{R}}
\providecommand{\Eta}{\textrm{H}}
\providecommand{\Beta}{\textrm{B}}
\providecommand{\Iota}{\textrm{J}}
\providecommand{\Omicron}{\textrm{O}}
\providecommand{\Kappa}{\textrm{K}}
\providecommand{\Mu}{\textrm{M}}
\providecommand{\Alpha}{\textrm{A}}
\providecommand{\Zeta}{\textrm{Z}}
\providecommand{\Tau}{\textrm{T}}
\providecommand{\Nu}{\textrm{N}}
\providecommand{\omicron}{\textrm{o}}


\begin{document}
\pagestyle{empty}\thispagestyle{empty}\lthtmltypeout{}%
\lthtmltypeout{latex2htmlLength hsize=\the\hsize}\lthtmltypeout{}%
\lthtmltypeout{latex2htmlLength vsize=\the\vsize}\lthtmltypeout{}%
\lthtmltypeout{latex2htmlLength hoffset=\the\hoffset}\lthtmltypeout{}%
\lthtmltypeout{latex2htmlLength voffset=\the\voffset}\lthtmltypeout{}%
\lthtmltypeout{latex2htmlLength topmargin=\the\topmargin}\lthtmltypeout{}%
\lthtmltypeout{latex2htmlLength topskip=\the\topskip}\lthtmltypeout{}%
\lthtmltypeout{latex2htmlLength headheight=\the\headheight}\lthtmltypeout{}%
\lthtmltypeout{latex2htmlLength headsep=\the\headsep}\lthtmltypeout{}%
\lthtmltypeout{latex2htmlLength parskip=\the\parskip}\lthtmltypeout{}%
\lthtmltypeout{latex2htmlLength oddsidemargin=\the\oddsidemargin}\lthtmltypeout{}%
\makeatletter
\if@twoside\lthtmltypeout{latex2htmlLength evensidemargin=\the\evensidemargin}%
\else\lthtmltypeout{latex2htmlLength evensidemargin=\the\oddsidemargin}\fi%
\lthtmltypeout{}%
\makeatother
\setcounter{page}{1}
\onecolumn

% !!! IMAGES START HERE !!!

\setcounter{secnumdepth}{4}
\setcounter{tocdepth}{5}
\stepcounter{paragraph}
\stepcounter{paragraph}
\stepcounter{section}
\stepcounter{section}
{\newpage\clearpage
\lthtmlfigureA{center3827}%
\begin{center}\vbox{\input{features_web.tex}
}\end{center}%
\lthtmlfigureZ
\lthtmlcheckvsize\clearpage}

\stepcounter{section}
\stepcounter{section}
\stepcounter{subsection}
\stepcounter{subsection}
\stepcounter{subsubsection}
\stepcounter{subsubsection}
\stepcounter{subsection}
{\newpage\clearpage
\lthtmlinlinemathA{tex2html_wrap_indisplay3836}%
$\displaystyle omi\  [options]\  [file]\  [args...]
$%
\lthtmlindisplaymathZ
\lthtmlcheckvsize\clearpage}

{\newpage\clearpage
\lthtmlinlinemathA{tex2html_wrap_indisplay3838}%
$\displaystyle /usr/local/bin/omi$%
\lthtmlindisplaymathZ
\lthtmlcheckvsize\clearpage}

{\newpage\clearpage
\lthtmlinlinemathA{tex2html_wrap_indisplay3840}%
$\displaystyle /usr/local/share/omi$%
\lthtmlindisplaymathZ
\lthtmlcheckvsize\clearpage}

\stepcounter{subsection}
\stepcounter{subsubsection}
\stepcounter{subsubsection}
\stepcounter{subsubsection}
\stepcounter{subsubsection}
\stepcounter{paragraph}
{\newpage\clearpage
\lthtmlinlinemathA{tex2html_wrap_inline3848}%
$ new$%
\lthtmlinlinemathZ
\lthtmlcheckvsize\clearpage}

\stepcounter{subsubsection}
\stepcounter{paragraph}
{\newpage\clearpage
\lthtmlinlinemathA{tex2html_wrap_inline2500}%
\fbox{\parbox[c]{\textwidth}{\tiny{
   Para más información:
   \begin{itemize}
   \setlength{\itemsep}{1px}
   \item Operadores lógicos (\autoref{sec:op_bool}).
   \item Operadores de comparación (\autoref{sec:op_cmp}).
   \item Conversión implícita entre tipos de datos (\autoref{sec:type_iconv})
   \end{itemize}
}}}%
\lthtmlinlinemathZ
\lthtmlcheckvsize\clearpage}

\stepcounter{paragraph}
\stepcounter{paragraph}
\stepcounter{paragraph}
{\newpage\clearpage
\lthtmlinlinemathA{tex2html_wrap_inline2546}%
\fbox{\parbox[c]{\textwidth}{\tiny{
   Para más información:
   \begin{itemize}
   \setlength{\itemsep}{1px}
   \item Operadores sobre arrays (\autoref{sec:op_string}).
   \item Funciones sobre arrays (\autoref{sec:func_string}).
   \end{itemize}
}}}%
\lthtmlinlinemathZ
\lthtmlcheckvsize\clearpage}

\stepcounter{paragraph}
{\newpage\clearpage
\lthtmlinlinemathA{tex2html_wrap_inline2560}%
\fbox{\parbox[c]{\textwidth}{\tiny{
   Para más información:
   \begin{itemize}
   \setlength{\itemsep}{1px}
   \item Clases de objetos (\autoref{sec:class}).
   \end{itemize}
}}}%
\lthtmlinlinemathZ
\lthtmlcheckvsize\clearpage}

\stepcounter{paragraph}
\stepcounter{paragraph}
{\newpage\clearpage
\lthtmlinlinemathA{tex2html_wrap_inline2580}%
\fbox{\parbox[c]{\textwidth}{\tiny{
   Para más información:
   \begin{itemize}
   \setlength{\itemsep}{1px}
   \item Funciones (\autoref{sec:function}).
   \end{itemize}
}}}%
\lthtmlinlinemathZ
\lthtmlcheckvsize\clearpage}

\stepcounter{paragraph}
{\newpage\clearpage
\lthtmlinlinemathA{tex2html_wrap_inline2592}%
\fbox{\parbox[c]{\textwidth}{\tiny{
   Para más información:
   \begin{itemize}
   \setlength{\itemsep}{1px}
   \item Funciones de expresiones regulares (\autoref{sec:func_regex}).
   \end{itemize}
}}}%
\lthtmlinlinemathZ
\lthtmlcheckvsize\clearpage}

\stepcounter{paragraph}
{\newpage\clearpage
\lthtmlinlinemathA{tex2html_wrap_inline3898}%
$ ''$%
\lthtmlinlinemathZ
\lthtmlcheckvsize\clearpage}

\stepcounter{subsubsection}
\stepcounter{paragraph}
{\newpage\clearpage
\lthtmlinlinemathA{tex2html_wrap_inline3982}%
$ !$%
\lthtmlinlinemathZ
\lthtmlcheckvsize\clearpage}

{\newpage\clearpage
\lthtmlinlinemathA{tex2html_wrap_inline3984}%
$ \&\&$%
\lthtmlinlinemathZ
\lthtmlcheckvsize\clearpage}

{\newpage\clearpage
\lthtmlinlinemathA{tex2html_wrap_inline3986}%
$ and$%
\lthtmlinlinemathZ
\lthtmlcheckvsize\clearpage}

{\newpage\clearpage
\lthtmlinlinemathA{tex2html_wrap_inline3988}%
$ ||$%
\lthtmlinlinemathZ
\lthtmlcheckvsize\clearpage}

{\newpage\clearpage
\lthtmlinlinemathA{tex2html_wrap_inline3990}%
$ or$%
\lthtmlinlinemathZ
\lthtmlcheckvsize\clearpage}

{\newpage\clearpage
\lthtmlinlinemathA{tex2html_wrap_inline2632}%
\fbox{\parbox[c]{\textwidth}{\tiny{
   Para más información:
   \begin{itemize}
   \setlength{\itemsep}{1px}
   \item Tipos de dato (\autoref{sec:data_type}).
   \item Tipo de dato booleano (\autoref{sec:type_bool}).
   \end{itemize}
}}}%
\lthtmlinlinemathZ
\lthtmlcheckvsize\clearpage}

\stepcounter{paragraph}
{\newpage\clearpage
\lthtmlinlinemathA{tex2html_wrap_inline3994}%
$ ==$%
\lthtmlinlinemathZ
\lthtmlcheckvsize\clearpage}

{\newpage\clearpage
\lthtmlinlinemathA{tex2html_wrap_inline3996}%
$ !=$%
\lthtmlinlinemathZ
\lthtmlcheckvsize\clearpage}

{\newpage\clearpage
\lthtmlinlinemathA{tex2html_wrap_inline3998}%
$ <$%
\lthtmlinlinemathZ
\lthtmlcheckvsize\clearpage}

{\newpage\clearpage
\lthtmlinlinemathA{tex2html_wrap_inline4000}%
$ <=$%
\lthtmlinlinemathZ
\lthtmlcheckvsize\clearpage}

{\newpage\clearpage
\lthtmlinlinemathA{tex2html_wrap_inline4002}%
$ >$%
\lthtmlinlinemathZ
\lthtmlcheckvsize\clearpage}

{\newpage\clearpage
\lthtmlinlinemathA{tex2html_wrap_inline4004}%
$ >=$%
\lthtmlinlinemathZ
\lthtmlcheckvsize\clearpage}

{\newpage\clearpage
\lthtmlinlinemathA{tex2html_wrap_inline4006}%
$ ===$%
\lthtmlinlinemathZ
\lthtmlcheckvsize\clearpage}

{\newpage\clearpage
\lthtmlinlinemathA{tex2html_wrap_inline4008}%
$ !==$%
\lthtmlinlinemathZ
\lthtmlcheckvsize\clearpage}

\stepcounter{paragraph}
{\newpage\clearpage
\lthtmlinlinemathA{tex2html_wrap_inline4012}%
$ +$%
\lthtmlinlinemathZ
\lthtmlcheckvsize\clearpage}

{\newpage\clearpage
\lthtmlinlinemathA{tex2html_wrap_inline4014}%
$ -$%
\lthtmlinlinemathZ
\lthtmlcheckvsize\clearpage}

{\newpage\clearpage
\lthtmlinlinemathA{tex2html_wrap_inline4016}%
$ *$%
\lthtmlinlinemathZ
\lthtmlcheckvsize\clearpage}

{\newpage\clearpage
\lthtmlinlinemathA{tex2html_wrap_inline4018}%
$ /$%
\lthtmlinlinemathZ
\lthtmlcheckvsize\clearpage}

{\newpage\clearpage
\lthtmlinlinemathA{tex2html_wrap_inline4020}%
$ \%$%
\lthtmlinlinemathZ
\lthtmlcheckvsize\clearpage}

{\newpage\clearpage
\lthtmlinlinemathA{tex2html_wrap_inline2670}%
\fbox{\parbox[c]{\textwidth}{\tiny{
   Para más información:
   \begin{itemize}
   \setlength{\itemsep}{1px}
   \item Tipos de dato (\autoref{sec:data_type}).
   \item Tipo de dato numérico (\autoref{sec:type_num}).
   \end{itemize}
}}}%
\lthtmlinlinemathZ
\lthtmlcheckvsize\clearpage}

\stepcounter{paragraph}
{\newpage\clearpage
\lthtmlinlinemathA{tex2html_wrap_inline4024}%
$ .$%
\lthtmlinlinemathZ
\lthtmlcheckvsize\clearpage}

{\newpage\clearpage
\lthtmlinlinemathA{tex2html_wrap_inline4026}%
$ [n]$%
\lthtmlinlinemathZ
\lthtmlcheckvsize\clearpage}

{\newpage\clearpage
\lthtmlinlinemathA{tex2html_wrap_inline2696}%
\fbox{\parbox[c]{\textwidth}{\tiny{
   Para más información:
   \begin{itemize}
   \setlength{\itemsep}{1px}
   \item Tipos de dato (\autoref{sec:data_type}).
   \item Tipo de dato cadena de caracteres (\autoref{sec:type_string}).
   \item Funciones sobre cadena de caracteres (\autoref{sec:fun_string}).
   \item Operadores sobre array (\autoref{sec:op_array}).
   \end{itemize}
}}}%
\lthtmlinlinemathZ
\lthtmlcheckvsize\clearpage}

\stepcounter{paragraph}
{\newpage\clearpage
\lthtmlinlinemathA{tex2html_wrap_inline4030}%
$ +=$%
\lthtmlinlinemathZ
\lthtmlcheckvsize\clearpage}

{\newpage\clearpage
\lthtmlinlinemathA{tex2html_wrap_inline4032}%
$ -=$%
\lthtmlinlinemathZ
\lthtmlcheckvsize\clearpage}

{\newpage\clearpage
\lthtmlinlinemathA{tex2html_wrap_inline4034}%
$ *=$%
\lthtmlinlinemathZ
\lthtmlcheckvsize\clearpage}

{\newpage\clearpage
\lthtmlinlinemathA{tex2html_wrap_inline4036}%
$ /=$%
\lthtmlinlinemathZ
\lthtmlcheckvsize\clearpage}

{\newpage\clearpage
\lthtmlinlinemathA{tex2html_wrap_inline4038}%
$ \%=$%
\lthtmlinlinemathZ
\lthtmlcheckvsize\clearpage}

{\newpage\clearpage
\lthtmlinlinemathA{tex2html_wrap_inline4040}%
$ .=$%
\lthtmlinlinemathZ
\lthtmlcheckvsize\clearpage}

{\newpage\clearpage
\lthtmlinlinemathA{tex2html_wrap_inline2712}%
\fbox{\parbox[c]{\textwidth}{\tiny{
   Para más información:
   \begin{itemize}
   \setlength{\itemsep}{1px}
   \item Operadores aritméticos (\autoref{sec:op_arth}).
   \item Operadores sobre cadenas de caracteres (\autoref{sec:op_string}).
   \end{itemize}
}}}%
\lthtmlinlinemathZ
\lthtmlcheckvsize\clearpage}

\stepcounter{paragraph}
{\newpage\clearpage
\lthtmlinlinemathA{tex2html_wrap_inline4044}%
$ ++i$%
\lthtmlinlinemathZ
\lthtmlcheckvsize\clearpage}

{\newpage\clearpage
\lthtmlinlinemathA{tex2html_wrap_inline4046}%
$ i++$%
\lthtmlinlinemathZ
\lthtmlcheckvsize\clearpage}

{\newpage\clearpage
\lthtmlinlinemathA{tex2html_wrap_inline4048}%
$ --i$%
\lthtmlinlinemathZ
\lthtmlcheckvsize\clearpage}

{\newpage\clearpage
\lthtmlinlinemathA{tex2html_wrap_inline4050}%
$ i--$%
\lthtmlinlinemathZ
\lthtmlcheckvsize\clearpage}

{\newpage\clearpage
\lthtmlinlinemathA{tex2html_wrap_inline2724}%
\fbox{\parbox[c]{\textwidth}{\tiny{
   Para más información:
   \begin{itemize}
   \setlength{\itemsep}{1px}
   \item Operadores aritméticos (\autoref{sec:op_arth}).
   \item Operadores de asignación (\autoref{sec:op_asig}).
   \end{itemize}
}}}%
\lthtmlinlinemathZ
\lthtmlcheckvsize\clearpage}

\stepcounter{paragraph}
{\newpage\clearpage
\lthtmlinlinemathA{tex2html_wrap_inline4054}%
$ bool$%
\lthtmlinlinemathZ
\lthtmlcheckvsize\clearpage}

{\newpage\clearpage
\lthtmlinlinemathA{tex2html_wrap_inline4056}%
$ int$%
\lthtmlinlinemathZ
\lthtmlcheckvsize\clearpage}

{\newpage\clearpage
\lthtmlinlinemathA{tex2html_wrap_inline4058}%
$ float$%
\lthtmlinlinemathZ
\lthtmlcheckvsize\clearpage}

{\newpage\clearpage
\lthtmlinlinemathA{tex2html_wrap_inline4060}%
$ string$%
\lthtmlinlinemathZ
\lthtmlcheckvsize\clearpage}

{\newpage\clearpage
\lthtmlinlinemathA{tex2html_wrap_inline2736}%
\fbox{\parbox[c]{\textwidth}{\tiny{
   Para más información:
   \begin{itemize}
   \setlength{\itemsep}{1px}
   \item Tipos de datos (\autoref{sec:data_type}).
   \item Conversión implícita entre tipos de datos (\autoref{sec:type_iconv}).
   \end{itemize}
}}}%
\lthtmlinlinemathZ
\lthtmlcheckvsize\clearpage}

\stepcounter{paragraph}
{\newpage\clearpage
\lthtmlinlinemathA{tex2html_wrap_inline4064}%
$ [k]$%
\lthtmlinlinemathZ
\lthtmlcheckvsize\clearpage}

{\newpage\clearpage
\lthtmlinlinemathA{tex2html_wrap_inline4066}%
$ []$%
\lthtmlinlinemathZ
\lthtmlcheckvsize\clearpage}

{\newpage\clearpage
\lthtmlinlinemathA{tex2html_wrap_inline4068}%
$ size$%
\lthtmlinlinemathZ
\lthtmlcheckvsize\clearpage}

{\newpage\clearpage
\lthtmlinlinemathA{tex2html_wrap_inline2760}%
\fbox{\parbox[c]{\textwidth}{\tiny{
   Para más información:
   \begin{itemize}
   \setlength{\itemsep}{1px}
   \item Tipos de dato (\autoref{sec:data_type}).
   \item Tipo de dato array (\autoref{sec:type_array}).
   \item Funciones sobre array (\autoref{sec:func_array}).
   \item Operadores sobre cadenas de caracteres (\autoref{sec:op_string}).
\par
\end{itemize}
}}}%
\lthtmlinlinemathZ
\lthtmlcheckvsize\clearpage}

\stepcounter{paragraph}
{\newpage\clearpage
\lthtmlinlinemathA{tex2html_wrap_inline4076}%
$ this$%
\lthtmlinlinemathZ
\lthtmlcheckvsize\clearpage}

{\newpage\clearpage
\lthtmlinlinemathA{tex2html_wrap_inline4078}%
$ parent$%
\lthtmlinlinemathZ
\lthtmlcheckvsize\clearpage}

{\newpage\clearpage
\lthtmlinlinemathA{tex2html_wrap_inline4080}%
$ static$%
\lthtmlinlinemathZ
\lthtmlcheckvsize\clearpage}

{\newpage\clearpage
\lthtmlinlinemathA{tex2html_wrap_inline2772}%
\fbox{\parbox[c]{\textwidth}{\tiny{
   Para más información:
   \begin{itemize}
   \setlength{\itemsep}{1px}
   \item Clases de objetos (\autoref{sec:class}).
   \item Tipos de datos (\autoref{sec:data_type}).
   \item Tipos de datos objetos (\autoref{sec:type_object}).
   \end{itemize}
}}}%
\lthtmlinlinemathZ
\lthtmlcheckvsize\clearpage}

\stepcounter{paragraph}
{\newpage\clearpage
\lthtmlinlinemathA{tex2html_wrap_inline4084}%
$ (p0,...)$%
\lthtmlinlinemathZ
\lthtmlcheckvsize\clearpage}

{\newpage\clearpage
\lthtmlinlinemathA{tex2html_wrap_inline4086}%
$ P[v1,...]$%
\lthtmlinlinemathZ
\lthtmlcheckvsize\clearpage}

{\newpage\clearpage
\lthtmlinlinemathA{tex2html_wrap_inline4088}%
$ \sim>$%
\lthtmlinlinemathZ
\lthtmlcheckvsize\clearpage}

{\newpage\clearpage
\lthtmlinlinemathA{tex2html_wrap_inline2784}%
\fbox{\parbox[c]{\textwidth}{\tiny{
   Para más información:
   \begin{itemize}
   \setlength{\itemsep}{1px}
   \item Funciones (\autoref{sec:function}).
   \item Tipos de datos (\autoref{sec:data_type}).
   \item Tipos de datos funciones (\autoref{sec:type_func}).
   \end{itemize}
}}}%
\lthtmlinlinemathZ
\lthtmlcheckvsize\clearpage}

\stepcounter{paragraph}
{\newpage\clearpage
\lthtmlinlinemathA{tex2html_wrap_inline4092}%
$ isnull$%
\lthtmlinlinemathZ
\lthtmlcheckvsize\clearpage}

{\newpage\clearpage
\lthtmlinlinemathA{tex2html_wrap_inline4094}%
$ isbool$%
\lthtmlinlinemathZ
\lthtmlcheckvsize\clearpage}

{\newpage\clearpage
\lthtmlinlinemathA{tex2html_wrap_inline4096}%
$ isnum$%
\lthtmlinlinemathZ
\lthtmlcheckvsize\clearpage}

{\newpage\clearpage
\lthtmlinlinemathA{tex2html_wrap_inline4098}%
$ isstring$%
\lthtmlinlinemathZ
\lthtmlcheckvsize\clearpage}

{\newpage\clearpage
\lthtmlinlinemathA{tex2html_wrap_inline4100}%
$ isarray$%
\lthtmlinlinemathZ
\lthtmlcheckvsize\clearpage}

{\newpage\clearpage
\lthtmlinlinemathA{tex2html_wrap_inline4102}%
$ isobject$%
\lthtmlinlinemathZ
\lthtmlcheckvsize\clearpage}

{\newpage\clearpage
\lthtmlinlinemathA{tex2html_wrap_inline2794}%
\fbox{\parbox[c]{\textwidth}{\tiny{
   Para más información:
   \begin{itemize}
   \setlength{\itemsep}{1px}
   \item Tipos de datos (\autoref{sec:data_type}).
   \end{itemize}
}}}%
\lthtmlinlinemathZ
\lthtmlcheckvsize\clearpage}

\stepcounter{paragraph}
{\newpage\clearpage
\lthtmlinlinemathA{tex2html_wrap_inline4106}%
$ cond\  ?\  op1:op2$%
\lthtmlinlinemathZ
\lthtmlcheckvsize\clearpage}

{\newpage\clearpage
\lthtmlinlinemathA{tex2html_wrap_inline4108}%
$ [[ op1, op2,...,opn ]]$%
\lthtmlinlinemathZ
\lthtmlcheckvsize\clearpage}

{\newpage\clearpage
\lthtmlinlinemathA{tex2html_wrap_inline2816}%
\fbox{\parbox[c]{\textwidth}{\tiny{
   Para más información:
   \begin{itemize}
   \setlength{\itemsep}{1px}
   \item Tipos de datos (\autoref{sec:data_type}).
   \item Operadores lógicos (\autoref{sec:op_bool}).
   \item Operadores de comparación (\autoref{sec:op_cmp}).
   \end{itemize}
}}}%
\lthtmlinlinemathZ
\lthtmlcheckvsize\clearpage}

\stepcounter{paragraph}
{\newpage\clearpage
\lthtmlinlinemathA{tex2html_wrap_inline4112}%
$ >>$%
\lthtmlinlinemathZ
\lthtmlcheckvsize\clearpage}

{\newpage\clearpage
\lthtmlinlinemathA{tex2html_wrap_inline4114}%
$ input$%
\lthtmlinlinemathZ
\lthtmlcheckvsize\clearpage}

{\newpage\clearpage
\lthtmlinlinemathA{tex2html_wrap_inline4116}%
$ <.$%
\lthtmlinlinemathZ
\lthtmlcheckvsize\clearpage}

{\newpage\clearpage
\lthtmlinlinemathA{tex2html_wrap_inline4118}%
$ <<$%
\lthtmlinlinemathZ
\lthtmlcheckvsize\clearpage}

{\newpage\clearpage
\lthtmlinlinemathA{tex2html_wrap_inline4120}%
$ echo$%
\lthtmlinlinemathZ
\lthtmlcheckvsize\clearpage}

\stepcounter{paragraph}
{\newpage\clearpage
\lthtmlinlinemathA{tex2html_wrap_inline4134}%
$ (p0,..)$%
\lthtmlinlinemathZ
\lthtmlcheckvsize\clearpage}

{\newpage\clearpage
\lthtmlinlinemathA{tex2html_wrap_inline4148}%
$ ++$%
\lthtmlinlinemathZ
\lthtmlcheckvsize\clearpage}

{\newpage\clearpage
\lthtmlinlinemathA{tex2html_wrap_inline4150}%
$ --$%
\lthtmlinlinemathZ
\lthtmlcheckvsize\clearpage}

{\newpage\clearpage
\lthtmlinlinemathA{tex2html_wrap_inline4160}%
$ [[op1,...]]$%
\lthtmlinlinemathZ
\lthtmlcheckvsize\clearpage}

{\newpage\clearpage
\lthtmlinlinemathA{tex2html_wrap_inline4162}%
$ ?:$%
\lthtmlinlinemathZ
\lthtmlcheckvsize\clearpage}

{\newpage\clearpage
\lthtmlinlinemathA{tex2html_wrap_inline4196}%
$ =$%
\lthtmlinlinemathZ
\lthtmlcheckvsize\clearpage}

\stepcounter{subsubsection}
{\newpage\clearpage
\lthtmlinlinemathA{tex2html_wrap_inline4211}%
$ goto$%
\lthtmlinlinemathZ
\lthtmlcheckvsize\clearpage}

{\newpage\clearpage
\lthtmlinlinemathA{tex2html_wrap_inline2862}%
\fbox{\parbox[c]{\textwidth}{\tiny{
   Para más información:
   \begin{itemize}
   \setlength{\itemsep}{1px}
   \item Sentencia $goto$\  (\autoref{sec:stmt_goto}).
   \end{itemize}
}}}%
\lthtmlinlinemathZ
\lthtmlcheckvsize\clearpage}

\stepcounter{subsubsection}
\stepcounter{paragraph}
{\newpage\clearpage
\lthtmlinlinemathA{tex2html_wrap_inline4218}%
$ if...\  else...$%
\lthtmlinlinemathZ
\lthtmlcheckvsize\clearpage}

{\newpage\clearpage
\lthtmlinlinemathA{tex2html_wrap_inline4222}%
$ else$%
\lthtmlinlinemathZ
\lthtmlcheckvsize\clearpage}

{\newpage\clearpage
\lthtmlinlinemathA{tex2html_wrap_inline4224}%
$ if...\  elif...\  else...$%
\lthtmlinlinemathZ
\lthtmlcheckvsize\clearpage}

{\newpage\clearpage
\lthtmlinlinemathA{tex2html_wrap_inline4226}%
$ elif$%
\lthtmlinlinemathZ
\lthtmlcheckvsize\clearpage}

\stepcounter{paragraph}
{\newpage\clearpage
\lthtmlinlinemathA{tex2html_wrap_inline4229}%
$ switch...\  case...$%
\lthtmlinlinemathZ
\lthtmlcheckvsize\clearpage}

{\newpage\clearpage
\lthtmlinlinemathA{tex2html_wrap_inline4235}%
$ break$%
\lthtmlinlinemathZ
\lthtmlcheckvsize\clearpage}

{\newpage\clearpage
\lthtmlinlinemathA{tex2html_wrap_inline4237}%
$ default$%
\lthtmlinlinemathZ
\lthtmlcheckvsize\clearpage}

{\newpage\clearpage
\lthtmlinlinemathA{tex2html_wrap_inline2902}%
\fbox{\parbox[c]{\textwidth}{\tiny{
   Para más información:
   \begin{itemize}
   \setlength{\itemsep}{1px}
   \item Sentencia break (\autoref{sec:stmt_break}).
   \end{itemize}
}}}%
\lthtmlinlinemathZ
\lthtmlcheckvsize\clearpage}

\stepcounter{paragraph}
{\newpage\clearpage
\lthtmlinlinemathA{tex2html_wrap_inline4241}%
$ while...$%
\lthtmlinlinemathZ
\lthtmlcheckvsize\clearpage}

\stepcounter{paragraph}
{\newpage\clearpage
\lthtmlinlinemathA{tex2html_wrap_inline4246}%
$ do...\  while$%
\lthtmlinlinemathZ
\lthtmlcheckvsize\clearpage}

\stepcounter{paragraph}
{\newpage\clearpage
\lthtmlinlinemathA{tex2html_wrap_inline4255}%
$ for...$%
\lthtmlinlinemathZ
\lthtmlcheckvsize\clearpage}

\stepcounter{paragraph}
{\newpage\clearpage
\lthtmlinlinemathA{tex2html_wrap_inline4264}%
$ for...\  as...$%
\lthtmlinlinemathZ
\lthtmlcheckvsize\clearpage}

{\newpage\clearpage
\lthtmlinlinemathA{tex2html_wrap_inline4266}%
$ for...\  in...$%
\lthtmlinlinemathZ
\lthtmlcheckvsize\clearpage}

\stepcounter{paragraph}
\stepcounter{paragraph}
\stepcounter{paragraph}
{\newpage\clearpage
\lthtmlinlinemathA{tex2html_wrap_inline4286}%
$ continue$%
\lthtmlinlinemathZ
\lthtmlcheckvsize\clearpage}

\stepcounter{paragraph}
{\newpage\clearpage
\lthtmlinlinemathA{tex2html_wrap_inline2992}%
\fbox{\parbox[c]{\textwidth}{\tiny{
   Para más información:
   \begin{itemize}
   \setlength{\itemsep}{1px}
   \item Etiquetas (\autoref{sec:label}).
   \item Sentencia condicional $if...\ else...$\  (\autoref{sec:stmt_if}).
   \end{itemize}
}}}%
\lthtmlinlinemathZ
\lthtmlcheckvsize\clearpage}

\stepcounter{paragraph}
{\newpage\clearpage
\lthtmlinlinemathA{tex2html_wrap_inline4303}%
$ include$%
\lthtmlinlinemathZ
\lthtmlcheckvsize\clearpage}

\stepcounter{paragraph}
{\newpage\clearpage
\lthtmlinlinemathA{tex2html_wrap_inline4308}%
$ exit$%
\lthtmlinlinemathZ
\lthtmlcheckvsize\clearpage}

\stepcounter{paragraph}
{\newpage\clearpage
\lthtmlinlinemathA{tex2html_wrap_inline4313}%
$ with$%
\lthtmlinlinemathZ
\lthtmlcheckvsize\clearpage}

\stepcounter{paragraph}
{\newpage\clearpage
\lthtmlinlinemathA{tex2html_wrap_inline4318}%
$ try...\  catch...$%
\lthtmlinlinemathZ
\lthtmlcheckvsize\clearpage}

{\newpage\clearpage
\lthtmlinlinemathA{tex2html_wrap_inline4322}%
$ catch$%
\lthtmlinlinemathZ
\lthtmlcheckvsize\clearpage}

{\newpage\clearpage
\lthtmlinlinemathA{tex2html_wrap_inline3028}%
\fbox{\parbox[c]{\textwidth}{\tiny{
   Para más información:
   \begin{itemize}
   \setlength{\itemsep}{1px}
   \item Excepciones (\autoref{sec:exception}).
   \item Sentencia  $throw$\  (\autoref{sec:stmt_throw}).
   \end{itemize}
}}}%
\lthtmlinlinemathZ
\lthtmlcheckvsize\clearpage}

\stepcounter{paragraph}
{\newpage\clearpage
\lthtmlinlinemathA{tex2html_wrap_inline4326}%
$ throw$%
\lthtmlinlinemathZ
\lthtmlcheckvsize\clearpage}

{\newpage\clearpage
\lthtmlinlinemathA{tex2html_wrap_inline4336}%
$ try...\  catch$%
\lthtmlinlinemathZ
\lthtmlcheckvsize\clearpage}

{\newpage\clearpage
\lthtmlinlinemathA{tex2html_wrap_inline3040}%
\fbox{\parbox[c]{\textwidth}{\tiny{
   Para más información:
   \begin{itemize}
   \setlength{\itemsep}{1px}
   \item Excepciones (\autoref{sec:exception}).
   \item Sentencia  $try...\ catch...$\  (\autoref{sec:stmt_try}).
   \end{itemize}
}}}%
\lthtmlinlinemathZ
\lthtmlcheckvsize\clearpage}

\stepcounter{paragraph}
{\newpage\clearpage
\lthtmlinlinemathA{tex2html_wrap_inline4340}%
$ sleep$%
\lthtmlinlinemathZ
\lthtmlcheckvsize\clearpage}

\stepcounter{paragraph}
{\newpage\clearpage
\lthtmlinlinemathA{tex2html_wrap_inline4345}%
$ typeof$%
\lthtmlinlinemathZ
\lthtmlcheckvsize\clearpage}

\stepcounter{paragraph}
{\newpage\clearpage
\lthtmlinlinemathA{tex2html_wrap_inline4350}%
$ datInfo$%
\lthtmlinlinemathZ
\lthtmlcheckvsize\clearpage}

{\newpage\clearpage
\lthtmlinlinemathA{tex2html_wrap_inline4356}%
$ datinfo$%
\lthtmlinlinemathZ
\lthtmlcheckvsize\clearpage}

\stepcounter{subsubsection}
{\newpage\clearpage
\lthtmlinlinemathA{tex2html_wrap_inline4359}%
$ global$%
\lthtmlinlinemathZ
\lthtmlcheckvsize\clearpage}

{\newpage\clearpage
\lthtmlinlinemathA{tex2html_wrap_inline3080}%
\fbox{\parbox[c]{\textwidth}{\tiny{
   Para más información:
   \begin{itemize}
   \setlength{\itemsep}{1px}
   \item Operador de asignación (\autoref{sec:op_asig}).
   \item Identificadores (\autoref{sec:id})
   \end{itemize}
}}}%
\lthtmlinlinemathZ
\lthtmlcheckvsize\clearpage}

\stepcounter{subsubsection}
{\newpage\clearpage
\lthtmlinlinemathA{tex2html_wrap_inline4363}%
$ \&$%
\lthtmlinlinemathZ
\lthtmlcheckvsize\clearpage}

\stepcounter{subsubsection}
{\newpage\clearpage
\lthtmlinlinemathA{tex2html_wrap_inline4368}%
$ \sim$%
\lthtmlinlinemathZ
\lthtmlcheckvsize\clearpage}

{\newpage\clearpage
\lthtmlinlinemathA{tex2html_wrap_inline4370}%
$ function$%
\lthtmlinlinemathZ
\lthtmlcheckvsize\clearpage}

{\newpage\clearpage
\lthtmlinlinemathA{tex2html_wrap_inline4372}%
$ return$%
\lthtmlinlinemathZ
\lthtmlcheckvsize\clearpage}

\stepcounter{paragraph}
\stepcounter{paragraph}
\stepcounter{paragraph}
\stepcounter{paragraph}
{\newpage\clearpage
\lthtmlinlinemathA{tex2html_wrap_inline4384}%
$ \sim\&$%
\lthtmlinlinemathZ
\lthtmlcheckvsize\clearpage}

\stepcounter{paragraph}
\stepcounter{paragraph}
\stepcounter{paragraph}
\stepcounter{paragraph}
\stepcounter{paragraph}
\stepcounter{paragraph}
{\newpage\clearpage
\lthtmlinlinemathA{tex2html_wrap_inline4400}%
$ \sim\sim$%
\lthtmlinlinemathZ
\lthtmlcheckvsize\clearpage}

\stepcounter{paragraph}
{\newpage\clearpage
\lthtmlinlinemathA{tex2html_wrap_inline4403}%
$ P$%
\lthtmlinlinemathZ
\lthtmlcheckvsize\clearpage}

\stepcounter{subsubsection}
{\newpage\clearpage
\lthtmlinlinemathA{tex2html_wrap_inline4408}%
$ for$%
\lthtmlinlinemathZ
\lthtmlcheckvsize\clearpage}

{\newpage\clearpage
\lthtmlinlinemathA{tex2html_wrap_inline4412}%
$ if$%
\lthtmlinlinemathZ
\lthtmlcheckvsize\clearpage}

{\newpage\clearpage
\lthtmlinlinemathA{tex2html_wrap_inline4414}%
$ :$%
\lthtmlinlinemathZ
\lthtmlcheckvsize\clearpage}

\stepcounter{subsubsection}
{\newpage\clearpage
\lthtmlinlinemathA{tex2html_wrap_inline4419}%
$ class$%
\lthtmlinlinemathZ
\lthtmlcheckvsize\clearpage}

{\newpage\clearpage
\lthtmlinlinemathA{tex2html_wrap_inline4423}%
$ ->$%
\lthtmlinlinemathZ
\lthtmlcheckvsize\clearpage}

\stepcounter{paragraph}
{\newpage\clearpage
\lthtmlinlinemathA{tex2html_wrap_inline4428}%
$ private$%
\lthtmlinlinemathZ
\lthtmlcheckvsize\clearpage}

\stepcounter{paragraph}
{\newpage\clearpage
\lthtmlinlinemathA{tex2html_wrap_inline4433}%
$ ::$%
\lthtmlinlinemathZ
\lthtmlcheckvsize\clearpage}

\stepcounter{paragraph}
{\newpage\clearpage
\lthtmlinlinemathA{tex2html_wrap_inline4436}%
$ extends$%
\lthtmlinlinemathZ
\lthtmlcheckvsize\clearpage}

\stepcounter{paragraph}
\stepcounter{paragraph}
\stepcounter{paragraph}
\stepcounter{paragraph}
\stepcounter{subsubsection}
{\newpage\clearpage
\lthtmlinlinemathA{tex2html_wrap_inline4456}%
$ getenv$%
\lthtmlinlinemathZ
\lthtmlcheckvsize\clearpage}

{\newpage\clearpage
\lthtmlinlinemathA{tex2html_wrap_inline4458}%
$ args$%
\lthtmlinlinemathZ
\lthtmlcheckvsize\clearpage}

\stepcounter{subsubsection}
{\newpage\clearpage
\lthtmlinlinemathA{tex2html_wrap_inline4463}%
$ try$%
\lthtmlinlinemathZ
\lthtmlcheckvsize\clearpage}

\stepcounter{subsubsection}
{\newpage\clearpage
\lthtmlinlinemathA{tex2html_wrap_inline4470}%
$ @$%
\lthtmlinlinemathZ
\lthtmlcheckvsize\clearpage}

{\newpage\clearpage
\lthtmlinlinemathA{tex2html_wrap_inline4472}%
$ eval$%
\lthtmlinlinemathZ
\lthtmlcheckvsize\clearpage}

\stepcounter{subsubsection}
{\newpage\clearpage
\lthtmlinlinemathA{tex2html_wrap_inline4475}%
$ getclass$%
\lthtmlinlinemathZ
\lthtmlcheckvsize\clearpage}

\stepcounter{subsubsection}
\stepcounter{paragraph}
\stepcounter{paragraph}
\stepcounter{paragraph}
\stepcounter{subsection}
\stepcounter{subsubsection}
\stepcounter{paragraph}
{\newpage\clearpage
\lthtmlfigureA{framed1042}%
\begin{framed}
\hfill \\$array$\  \textbf{str\_explode} ( $string$\  \textbf{str}, $string$\  \textbf{delimeter})  
\begin{description}
\item [Parámetros:] \hfill 
   \begin{description}
   \item[str:] Cadena de caracteres de entrada.
   \item[delimeter:] Cadena delimitadora.
   \end{description}
\item[Valores devueltos:] \hfill \\
   Array de elementos que contiene las subcadenas resultado de dividir \textbf{str} usando el delimitador \textbf{delimeter}.
\par
Si el separador no se encuentra en la cadena el resultado será un array de un elemento que contendrá la cadena completa.
\end{description}
\end{framed}%
\lthtmlfigureZ
\lthtmlcheckvsize\clearpage}

\stepcounter{paragraph}
{\newpage\clearpage
\lthtmlfigureA{framed1056}%
\begin{framed}
\hfill \\$numeric$\  \textbf{str\_find} ( $string$\  \textbf{str}, $mixed$\  \textbf{substr} [, $numeric$\  \textbf{position}] )  
\begin{description}
\item [Parámetros:] \hfill 
   \begin{description}
   \item[str:] Cadena de caracteres de entrada.
   \item[substr:] Expresión de búsqueda. Puede ser una cadena de caracteres o una expresión regular.  
   \item[position:] Posición dentro de la cadena \textbf{str} a partir de la cual se comenzará la búsqueda. Si no es dado 
   se comenzará desde la primera posición de la cadena.
   \end{description}
\item[Valores devueltos:] \hfill \\
   Devuelve el número entero correspondiente a la primera posición en la que se encuentre \textbf{substr} dentro de \textbf{str}. 
   Se ha de tener en cuenta que la primera posición de la cadena se corresponde con el entero $0$.
\par
Si \textbf{substr} es una expresión regular esta se corresponderá con la subcadena de mayor longitud perteneciente al 
   lenguaje definido por esta. 
\par
Si la subcadena no se encuentra devuelve el valor $-1$. 
\par
Si la posición de inicio \textbf{position} es mayor que la longitud de la cadena \textbf{str} se devuelve el valor $-1$. 
\end{description}
\end{framed}%
\lthtmlfigureZ
\lthtmlcheckvsize\clearpage}

\stepcounter{paragraph}
{\newpage\clearpage
\lthtmlfigureA{framed1075}%
\begin{framed}
\hfill \\$string$\  \textbf{str\_replace} ( $string$\  \textbf{str}, $mixed$\  \textbf{search}, $string$\  \textbf{replace} [, $numeric$\  \textbf{n}] )  
\begin{description}
\item [Parámetros:] \hfill 
   \begin{description}
   \item[str:] Cadena de caracteres de entrada.
   \item[search:] Expresión de búsqueda. Puede ser una cadena de caracteres o una expresión regular.  
   \item[replace:] Cadena de remplazo.
   \item[n:] Número de reemplazos a efectuar. Si no es dado se realizan todos los reemplazos.
   \end{description}
\item[Valores devueltos:] \hfill \\
   Devuelve un cadena de caracteres resultado de remplazar la subcadena \textbf{search}  por \textbf{replace}
   en la cadena de caracteres \textbf{str}, tantas veces como indique el parámetro \textbf{n}.
\par
Si \textbf{search} es una expresión regular se buscarán subcadenas que pertenezcan al conjunto de palabras del lenguaje definido por esta. La
   expresión regular puede estar compuesta por subexpresiones delimitadas por paréntesis, si es así es posible hacer referencia a las subcadenas 
   correspondientes a estas usando el carácter ``\textbackslash'' seguido de la posición que ocupa la subexpresión.
\par
\end{description}
\end{framed}%
\lthtmlfigureZ
\lthtmlcheckvsize\clearpage}

\stepcounter{paragraph}
{\newpage\clearpage
\lthtmlfigureA{framed1094}%
\begin{framed}
\hfill \\$string$\  \textbf{str\_replace\_sub} ( $string$\  \textbf{str}, $numeric$\  \textbf{ini}, $numeric$\  \textbf{len}, $string$\  \textbf{replace} )  
\begin{description}
\item [Parámetros:] \hfill 
   \begin{description}
   \item[str:] Cadena de caracteres de entrada.
   \item[ini:] Posición inicial para el remplazo.
   \item[len:] Longitud de la subcadena de remplazo.
   \item[replace:] Cadena de remplazo.
   \end{description}
\item[Valores devueltos:] \hfill \\
   Devuelve un cadena de caracteres resultado de remplazar la subcadena dada por la posición \textbf{ini} y la longitud
   \textbf{len} dentro de la cadena de caracteres \textbf{str} por la cadena \textbf{replace}.
\par
Si la posición inicial dada se encuentra fuera de la cadena se produce un error semántico.
\end{description}
\end{framed}%
\lthtmlfigureZ
\lthtmlcheckvsize\clearpage}

\stepcounter{paragraph}
{\newpage\clearpage
\lthtmlfigureA{framed1112}%
\begin{framed}
\hfill \\$string$\  \textbf{str\_upper} ( $string$\  \textbf{str} )  
\begin{description}
\item [Parámetros:] \hfill 
   \begin{description}
   \item[str:] Cadena de caracteres de entrada.
   \end{description}
\item[Valores devueltos:] \hfill \\
   Devuelve la cadena de caracteres \textbf{str} con todos los caracteres alfabéticos convertidos a mayúsculas.
\end{description}
\end{framed}%
\lthtmlfigureZ
\lthtmlcheckvsize\clearpage}

\stepcounter{paragraph}
{\newpage\clearpage
\lthtmlfigureA{framed1124}%
\begin{framed}
\hfill \\$string$\  \textbf{str\_lower} ( $string$\  \textbf{str} )  
\begin{description}
\item [Parámetros:] \hfill 
   \begin{description}
   \item[str:] Cadena de caracteres de entrada.
   \end{description}
\item[Valores devueltos:] \hfill \\
   Devuelve la cadena de caracteres \textbf{str} con todos los caracteres alfabéticos convertidos a minúsculas.
\end{description}
\end{framed}%
\lthtmlfigureZ
\lthtmlcheckvsize\clearpage}

\stepcounter{paragraph}
{\newpage\clearpage
\lthtmlfigureA{framed1136}%
\begin{framed}
\hfill \\$array$\  \textbf{str\_search} ( $mixed$\  \textbf{str}, $regexp$\  \textbf{pattern} [, $string$\  \textbf{keys...}] )  
\begin{description}
\item [Parámetros:] \hfill 
   \begin{description}
   \item[str:] Cadena de entrada. Puede ser un array de cadenas.
   \item[pattern:] Patrón de búsqueda. 
   \item[keys...:] Listado de claves para las coincidencias. Si no se da las
   claves serán numéricas.
   \end{description}
\item[Valores devueltos:] \hfill \\
   Devuelve un array que contiene todas las coincidencias del patrón dado por la 
   expresión regular \textbf{pattern} en la cadena \textbf{str}, donde las 
   claves del array vienen dadas por la lista de cadenas \textbf{keys...}. \\
\par
La expresión regular puede estar formada por subexpresiones delimitadas
   por "()". En dicho caso se buscará en la cadena \textbf{str} subcadenas que pertenezcan al
   conjunto delimitado por la expresión regular. Por cada subcadena encontrada se creará un array
   con las correspondencias de cada subexpresión. \\
\par
Si \textbf{str} es un array de cadenas se aplicará el algoritmo de búsqueda de foram iterativa a cada cadena en el mismo.
\end{description}
\end{framed}%
\lthtmlfigureZ
\lthtmlcheckvsize\clearpage}

\stepcounter{paragraph}
{\newpage\clearpage
\lthtmlfigureA{framed1154}%
\begin{framed}
\hfill \\$bool$\  \textbf{str\_match} ( $string$\  \textbf{str}, $regexp$\  \textbf{pattern})  
\begin{description}
\item [Parámetros:] \hfill 
   \begin{description}
   \item[str:] Cadena de entrada.
   \item[pattern:] Expresión regular patrón. 
   \end{description}
\item[Valores devueltos:] \hfill \\
   Devuelve un valor booleano verdadero si la cadena de caracteres \textbf{str} pertenence
   al conjunto de palabras del lenguaje definido por la expresión regular \textbf{pattern}. Falso
   en caso contrario.
\end{description}
\end{framed}%
\lthtmlfigureZ
\lthtmlcheckvsize\clearpage}

\stepcounter{paragraph}
{\newpage\clearpage
\lthtmlfigureA{framed1168}%
\begin{framed}
\hfill \\$regexp$\  \textbf{regexp} ( $string$\  \textbf{str})  
\begin{description}
\item [Parámetros:] \hfill 
   \begin{description}
   \item[str:] Cadena de entrada.
   \end{description}
\item[Valores devueltos:] \hfill \\
   Devuelve la expresión regular correspondiente a la cadena \textbf{str}.
\end{description}
\end{framed}%
\lthtmlfigureZ
\lthtmlcheckvsize\clearpage}

\stepcounter{paragraph}
{\newpage\clearpage
\lthtmlfigureA{framed1180}%
\begin{framed}
\hfill \\$string$\  \textbf{sprintf} ( $string$\  \textbf{format} [, $mixed$\  \textbf{values...}] )  
\begin{description}
\item [Parámetros:] \hfill 
   \begin{description}
   \item[format:]  \hfill \\
      La cadena formato contendrá una serie de
		directivas de formato. Estas directivas serán sustituidas por el valor correspondiente, según posición, de la
		lista. Cuando se realiza cada sustitución el valor es formateado según la directiva. \\
\par
Las directivas de formato tienen el siguiente forma:
		\begin{displaymath}\%[operador][precisi\acute{o}n][formato]\end{displaymath}\  \\
		Los posibles operadores serán los siguientes:
		\begin{description}
			\item[ +:] Fuerza la impresión del símbolo + cuando se formatean números positivos.
			\item[ \^\ :] Convierte el caracteres a mayúsculas cuando se formatean cadenas de texto.
			\item[ \#:] Añade el carácter 0x cuando se formatean números hexadecimales y el carácter 0 cuando se formatean octales.
		\end{description} \hfill 
\par
La precisión se refiere al número de decimales que se imprimirán en el caso de formatear
		números o el número de caracteres en el caso de formatear cadenas. \\
\par
El carácter de formato indica que tipo de formato se le dará al valor: 
		\begin{description}
			\item[  i|d:] Número entero.
			\item[  u:] Sin signo.
			\item[  f:] Coma flotante.
			\item[  \%:] Carácter \%.
			\item[  e:] Notación científica.
			\item[  o:] Octal.
			\item[  x:] Hexadecimal.
			\item[  s|c:] Cadena de texto.
		\end{description} \hfill \\
   \item[values:] \hfill \\Lista de valores. Se deben de dar tantos como directivas existan en la cadena de formato.
\par
\end{description}
\item[Valores devueltos:] \hfill \\
   Devuelve la cadena de caracteres resultante de sustituir las directivas de formato de la cadena \textbf{format}, por los valores  de la lista \textbf{values} formateados 
   según la directiva correspondiente posicionalmente.
\end{description}
\end{framed}%
\lthtmlfigureZ
\lthtmlcheckvsize\clearpage}

\stepcounter{subsubsection}
\stepcounter{paragraph}
{\newpage\clearpage
\lthtmlfigureA{framed1200}%
\begin{framed}
\hfill \\$string$\  \textbf{array\_implode} ( $array$\  \textbf{elements}, $string$\  \textbf{sep} )  
\begin{description}
\item [Parámetros:] \hfill 
   \begin{description}
   \item[elements:] Array de elementos
   \item[sep:] Cadena separadora
   \end{description}
\item[Valores devueltos:] \hfill \\
   Devuelve un string correspondiente a todos los elementos del array \textbf{elements} en el mismo orden, con el string \textbf{sep} entre cada elemento. 
\end{description}
\end{framed}%
\lthtmlfigureZ
\lthtmlcheckvsize\clearpage}

\stepcounter{paragraph}
{\newpage\clearpage
\lthtmlfigureA{framed1214}%
\begin{framed}
\hfill \\$mixed$\  \textbf{array\_first} ( $array$\  \textbf{elements} )  
\begin{description}
\item [Parámetros:] \hfill 
   \begin{description}
   \item[elements:] Array de elementos
   \end{description}
\item[Valores devueltos:] \hfill \\
   Devuelve el primer elemento del array \textbf{elements}. Si \textbf{elements} es un array
   numérico con claves numéricas devuelve el elemento con la posición 0. Si es un array asociativo 
   devuelve el elemento correspondiente a la primera clave que se introdujo. Si el array es vacío 
   devuelve un valor nulo. 
\end{description}
\end{framed}%
\lthtmlfigureZ
\lthtmlcheckvsize\clearpage}

\stepcounter{paragraph}
{\newpage\clearpage
\lthtmlfigureA{framed1227}%
\begin{framed}
\hfill \\$mixed$\  \textbf{array\_last} ( $array$\  \textbf{elements} )  
\begin{description}
\item [Parámetros:] \hfill 
   \begin{description}
   \item[elements:] Array de elementos
   \end{description}
\item[Valores devueltos:] \hfill \\
   Devuelve el último elemento del array \textbf{elements}. Si \textbf{elements} es un array
   numérico con claves numéricas devuelve el elemento con la posición $size - 1$. Si es un array asociativo 
   devuelve el elemento correspondiente a la última clave que se introdujo. Si el array es vacío 
   devuelve un valor nulo. 
\end{description}
\end{framed}%
\lthtmlfigureZ
\lthtmlcheckvsize\clearpage}

\stepcounter{paragraph}
{\newpage\clearpage
\lthtmlfigureA{framed1240}%
\begin{framed}
\hfill \\$array$\  \textbf{array\_insert} ( $array$\  \&\textbf{elements}, $numeric$\  \textbf{position}, $mixed$\  \textbf{element})  
\begin{description}
\item [Parámetros:] \hfill 
   \begin{description}
   \item[elements:] Array de elementos. Es una referencia por lo que el array dado como valor de este parámetro será modificado. 
   \item[position:] Entero sin signo que determina la posición en la que se insertará el elemento. Debe ser un valor entero entre $0$\  y
   el tamaño del array, cualquier otro valor dará error de índice.
   \item[element:] Elemento a insertar
   \end{description}
\item[Valores devueltos:] \hfill \\
   Devuelve el array resultado de insertar \textbf{element} en la posición \textbf{position} del array \textbf{elements}.
   Si \textbf{elements} es un array de índices numéricos inserta el elemento en la posición dada. Si 
   es un array asociativo inserta el elemento en la clave que ocupa la posición dada, desplazando los demás valores y creando
   una nueva clave correspondiente al número de elementos.
\end{description}
\end{framed}%
\lthtmlfigureZ
\lthtmlcheckvsize\clearpage}

\stepcounter{paragraph}
{\newpage\clearpage
\lthtmlfigureA{framed1257}%
\begin{framed}
\hfill \\$array$\  \textbf{array\_delete} ( $array$\  \&\textbf{elements}, $mixed$\  \textbf{position})  
\begin{description}
\item [Parámetros:] \hfill 
   \begin{description}
   \item[elements:] Array de elementos. Es una referencia por lo que el array dado como valor de este parámetro será modificado. 
   \item[position:] Puede ser un entero sin signo que determina la posición que será eliminada. Debe ser un valor entero mayor que $0$\  y
   menor que el tamaño del array, cualqueir otro valor no tendrá efecto. Además puede ser una cadena de caracteres que 
   represente la clave del array que será eliminada. Si la clave no existe en el array la función no realizará ninguna acción.
   \end{description}
\item[Valores devueltos:] \hfill \\
   Si el array \textbf{elements} presenta índices numéricos devuelve el array resultado de eliminar la posición \textbf{position} del array. 
\par
Si el array \textbf{elements} es asociativo devuelve el array resultado de eliminar la clave \textbf{position} del array. 
\end{description}
\end{framed}%
\lthtmlfigureZ
\lthtmlcheckvsize\clearpage}

\stepcounter{paragraph}
{\newpage\clearpage
\lthtmlfigureA{framed1273}%
\begin{framed}
\hfill \\$array$\  \textbf{array\_push} ( $array$\  \&\textbf{elements}, $mixed$\  \textbf{element})  
\begin{description}
\item [Parámetros:] \hfill 
   \begin{description}
   \item[elements:] Array de elementos. Es una referencia por lo que el array dado como valor de este parámetro será modificado. 
   \item[element:] Elemento a insertar al final del array
   \end{description}
\item[Valores devueltos:] \hfill \\
   Devuelve el array \textbf{elements} con el elemento \textbf{element} insertado en la última posición. Si \textbf{elements} es
   asociativo la clave del elemento insertado será el número de elementos del array antes de la insercción. 
\end{description}
\end{framed}%
\lthtmlfigureZ
\lthtmlcheckvsize\clearpage}

\stepcounter{paragraph}
{\newpage\clearpage
\lthtmlfigureA{framed1288}%
\begin{framed}
\hfill \\$mixed$\  \textbf{array\_pop} ( $array$\  \&\textbf{elements} )  
\begin{description}
\item [Parámetros:] \hfill 
   \begin{description}
   \item[elements:] Array de elementos. Es una referencia por lo que el array dado como valor de este parámetro será modificado. 
   \end{description}
\item[Valores devueltos:] \hfill \\
   Devuelve el último elemento del array \textbf{elements} y lo elimina del mismo. Si el array es asociativo devuelve el
   elemento de la última clave introducida.
\end{description}
\end{framed}%
\lthtmlfigureZ
\lthtmlcheckvsize\clearpage}

\stepcounter{paragraph}
{\newpage\clearpage
\lthtmlfigureA{framed1300}%
\begin{framed}
\hfill \\$array$\  \textbf{array\_unshift} ( $array$\  \&\textbf{elements}, $mixed$\  \textbf{element})  
\begin{description}
\item [Parámetros:] \hfill 
   \begin{description}
   \item[elements:] Array de elementos. Es una referencia por lo que el array dado como valor de este parámetro será modificado. 
   \item[element:] Elemento a insertar al inicio del array
   \end{description}
\item[Valores devueltos:] \hfill \\
   Devuelve el array \textbf{elements} con el elemento \textbf{element} insertado en la primera posición. Si \textbf{elements} es
   asociativo la clave del elemento insertado será la misma que la anterior, llevándose a cabo un desplazamiento. 
\end{description}
\end{framed}%
\lthtmlfigureZ
\lthtmlcheckvsize\clearpage}

\stepcounter{paragraph}
{\newpage\clearpage
\lthtmlfigureA{framed1315}%
\begin{framed}
\hfill \\$mixed$\  \textbf{array\_shift} ( $array$\  \&\textbf{elements} )  
\begin{description}
\item [Parámetros:] \hfill 
   \begin{description}
   \item[elements:] Array de elementos. Es una referencia por lo que el array dado como valor de este parámetro será modificado. 
   \end{description}
\item[Valores devueltos:] \hfill \\
   Devuelve el primer elemento del array \textbf{elements} y lo elimina del mismo. Si el array es asociativo devuelve el
   elemento de la primera clave introducida.
\end{description}
\end{framed}%
\lthtmlfigureZ
\lthtmlcheckvsize\clearpage}

\stepcounter{paragraph}
{\newpage\clearpage
\lthtmlfigureA{framed1327}%
\begin{framed}
\hfill \\$mixed$\  \textbf{array\_reduce} ( $array$\  \textbf{elements}, $function$\  \textbf{iterative} )  
\begin{description}
\item [Parámetros:] \hfill 
   \begin{description}
   \item[elements:] Array de elementos.
   \item[itereative:] Función de iteración para la reducción. La función debe recibir dos parámetros donde el primero
   será el acumulador y el segundo el elemento de la iteración actual, y devolverá el acumulador para la próxima iteración.
   La iteración comienza con el primer elemento como valor del acumulador y el segundo como primer elemento para iterar.
   \end{description}
\item[Valores devueltos:] \hfill \\
   Devuelve el resultado de aplicar la función \textbf{iterative} iterativamente sobre los elementos del array \textbf{elements}.
\end{description}
\end{framed}%
\lthtmlfigureZ
\lthtmlcheckvsize\clearpage}

\stepcounter{subsubsection}
\stepcounter{paragraph}
{\newpage\clearpage
\lthtmlfigureA{framed1342}%
\begin{framed}
\hfill \\$string$\  \textbf{exec} ( $string$\  \textbf{cmd} )  
\begin{description}
\item [Parámetros:] \hfill 
   \begin{description}
   \item[cmd:] Cadena de caracteres que representa el comando a ejecutar
   \end{description}
\item[Valores devueltos:] \hfill \\
   Devuelve el contenido de la salida estándar al ejecutar el comando del sistema indicado por \textbf{cmd}.
\end{description}
\end{framed}%
\lthtmlfigureZ
\lthtmlcheckvsize\clearpage}

\stepcounter{paragraph}
{\newpage\clearpage
\lthtmlfigureA{framed1354}%
\begin{framed}
\hfill \\$numeric$\  \textbf{fork} ( )  
\begin{description}
\item[Valores devueltos:] \hfill \\
   Produce una bifurcación en la ejecución, creando un nuevo proceso copia del 
   proceso actual. La ejecución en el proceso hijo prosigue en la misma sentencia 
   y con el mismo estado que el proceso padre. 
\par
El valor devuelto en el proceso padre será de el identificador del proceso hijo creado. 
\par
El valor devuelto en el proceso hijo será $0$.
\end{description}
\end{framed}%
\lthtmlfigureZ
\lthtmlcheckvsize\clearpage}

\stepcounter{paragraph}
{\newpage\clearpage
\lthtmlfigureA{framed1362}%
\begin{framed}
\hfill \\$numeric$\  \textbf{wait} ( [ $numeric$\  pid ] )  
\begin{description}
\item [Parámetros:] \hfill 
   \begin{description}
   \item[pid:] Identificador del proceso hijo cuya ejecución se desea esperar a que 
   finalice. Si no es dado se esperará a que finalicen todos los procesos hijos.
   \end{description}
\item[Valores devueltos:] \hfill \\
   Devuelve el código de salida producido por el último proceso en finalizar.
\end{description}
\end{framed}%
\lthtmlfigureZ
\lthtmlcheckvsize\clearpage}

\stepcounter{paragraph}
{\newpage\clearpage
\lthtmlfigureA{framed1372}%
\begin{framed}
\hfill \\$numeric$\  \textbf{getpid} ( )  
\begin{description}
\item[Valores devueltos:] \hfill \\
   Devuelve el identificador del proceso actual
\end{description}
\end{framed}%
\lthtmlfigureZ
\lthtmlcheckvsize\clearpage}

\stepcounter{paragraph}
{\newpage\clearpage
\lthtmlfigureA{framed1380}%
\begin{framed}
\hfill \\$numeric$\  \textbf{getppid} ( )  
\begin{description}
\item[Valores devueltos:] \hfill \\
   Devuelve el identificador del proceso padre del actual.
\end{description}
\end{framed}%
\lthtmlfigureZ
\lthtmlcheckvsize\clearpage}

\stepcounter{paragraph}
{\newpage\clearpage
\lthtmlfigureA{framed1388}%
\begin{framed}
\hfill \\$bool$\  \textbf{signal} ( $numeric$\  \textbf{pid}, $numeric$\  \textbf{code} )  
\begin{description}
\item [Parámetros:] \hfill 
   \begin{description}
   \item[pid:] Identificador de proceso al que se le envía la señal
   \item[code:] Código numérico de señal POSIX
   \end{description}
\item[Valores devueltos:] \hfill \\
   Un valor booleano que indica si la señal se envió correctamente. 
\end{description}
\end{framed}%
\lthtmlfigureZ
\lthtmlcheckvsize\clearpage}

\stepcounter{paragraph}
{\newpage\clearpage
\lthtmlfigureA{framed1400}%
\begin{framed}
\hfill \\$bool$\  \textbf{shandler} ( $numeric$\  \textbf{code}, $function$\  \textbf{handler} )  
\begin{description}
\item [Parámetros:] \hfill 
   \begin{description}
   \item[code:] Código numérico de señal POSIX a la que se le atribuirá el manejador.
   \item[handler:] Función que manejará la ejecución cuando se dé la señal.
   \end{description}
\item[Valores devueltos:] \hfill \\
   Un valor booleano que indica si el manejador de señal se estableció correctamente. 
\end{description}
\end{framed}%
\lthtmlfigureZ
\lthtmlcheckvsize\clearpage}

\stepcounter{paragraph}
{\newpage\clearpage
\lthtmlfigureA{framed1412}%
\begin{framed}
\hfill \\$void$\  \textbf{exit\_process} ( )  
\end{framed}%
\lthtmlfigureZ
\lthtmlcheckvsize\clearpage}

\stepcounter{paragraph}
{\newpage\clearpage
\lthtmlfigureA{framed1418}%
\begin{framed}
\hfill \\$bool$\  \textbf{process} ( $function$\  \textbf{func} [, $mixed$\  \textbf{args ...} ]  )  
\begin{description}
\item [Parámetros:] \hfill 
   \begin{description}
   \item[func:] Función que será llamada como un proceso.
   \item[args...:] Lista de parámetros que serán pasados a la función.
   \end{description}
\item[Valores devueltos:] \hfill \\
   Devuelve un valor verdadero si el proceso fue llamado con éxito.
\end{description}
\end{framed}%
\lthtmlfigureZ
\lthtmlcheckvsize\clearpage}

\stepcounter{subsubsection}
\stepcounter{paragraph}
{\newpage\clearpage
\lthtmlfigureA{framed1431}%
\begin{framed}
\hfill \\$file$\  \textbf{file} ( $string$\  \textbf{name}, $string$\  \textbf{mode}  )  
\begin{description}
\item [Parámetros:] \hfill 
   \begin{description}
   \item[name:] Nombre del fichero que se desea abrir. Se puede utilizar una ruta absoluta o relativa al programa
   \item[mode:] Modo en el que se abrirá el fichero:
      \begin{description}
         \item [r:] Solo lectura
         \item [r+:] Lectura y/o escritura
         \item [w:] Escritura truncando el fichero
         \item [w+:] Lectura y/o escritura truncando el fichero
         \item [a:] Escritura posicionando el puntero al final del fichero
         \item [a+:] Lectura y/o escritura posicionando el puntero al final del fichero
      \end{description} \hfill 
\par
Todos los modos excepto el solo lectura crean el fichero si este no existe. 
   \end{description}
\item[Valores devueltos:] \hfill \\
   Devuelve un puntero a fichero si el fichero ha sido abierto correctamente. Y un valor falso si no se ha 
   abierto el fichero con éxito
\end{description}
\end{framed}%
\lthtmlfigureZ
\lthtmlcheckvsize\clearpage}

\stepcounter{paragraph}
{\newpage\clearpage
\lthtmlfigureA{framed1445}%
\begin{framed}
\hfill \\$bool$\  \textbf{fclose} ( $file$\  \textbf{fpointer}  )  
\begin{description}
\item [Parámetros:] \hfill 
   \begin{description}
   \item[fpointer:] Puntero a fichero. 
   \end{description}
\item[Valores devueltos:] \hfill \\
   Devuelve un valor booleano que indica si el fichero se ha cerrado correctamente.
\end{description}
\end{framed}%
\lthtmlfigureZ
\lthtmlcheckvsize\clearpage}

\stepcounter{paragraph}
{\newpage\clearpage
\lthtmlfigureA{framed1456}%
\begin{framed}
\hfill \\$numeric$\  \textbf{fput}|\textbf{f>:} ( $file$\  \textbf{fpointer}, $string$\  \textbf{content}  )  
\begin{description}
\item [Parámetros:] \hfill 
   \begin{description}
   \item[fpointer:] Puntero a fichero. 
   \item[content:] Cadena de caracteres que será escrita en el fichero a partir de la posición en la que se encuentre el puntero. 
   \end{description}
\item[Valores devueltos:] \hfill \\
   Devuelve un valor numérico que representa los bytes que han sido escritos.
\end{description}
\end{framed}%
\lthtmlfigureZ
\lthtmlcheckvsize\clearpage}

\stepcounter{paragraph}
{\newpage\clearpage
\lthtmlfigureA{framed1469}%
\begin{framed}
\hfill \\$string$\  \textbf{fget}|\textbf{f<:} ( $file$\  \textbf{fpointer} [, $numeric$\  \textbf{size} ]  )  
\begin{description}
\item [Parámetros:] \hfill 
   \begin{description}
   \item[fpointer:] Puntero a fichero. 
   \item[size:] Número de caracteres que serán leídos. Si no es dado se lee hasta el primer carácter de salto de línea encontrado o hasta el final del fichero 
   \end{description}
\item[Valores devueltos:] \hfill \\
   Devuelve la cadena de caracteres leída.
\end{description}
\end{framed}%
\lthtmlfigureZ
\lthtmlcheckvsize\clearpage}

\stepcounter{paragraph}
{\newpage\clearpage
\lthtmlfigureA{framed1482}%
\begin{framed}
\hfill \\$bool$\  \textbf{fseek} ( $file$\  \textbf{fpointer}, $numeric$\  \textbf{offset} [, $enum$\  \textbf{position} ]  )  
\begin{description}
\item [Parámetros:] \hfill 
   \begin{description}
   \item[fpointer:] Puntero a fichero. 
   \item[offset:] Número de bytes que será desplazado el puntero
   \item[position:] Posición de referencia desde la que mover el puntero. Puede ser algún valor 
   FSET, FCUR o FEND, para referirse al inicio, la posición actual o el final. Si no es dado se
   toma como referencia el inicio del fichero.
   \end{description}
\item[Valores devueltos:] \hfill \\
   Devuelve un booleano que indica si la posición cambió correctamente. Si la posición 
   nueva se encuentra fuera del fichero el puntero se colocará al final del mismo.
\end{description}
\end{framed}%
\lthtmlfigureZ
\lthtmlcheckvsize\clearpage}

\stepcounter{paragraph}
{\newpage\clearpage
\lthtmlfigureA{framed1495}%
\begin{framed}
\hfill \\$numeric$\  \textbf{ftell} ( $file$\  \textbf{fpointer}  )  
\begin{description}
\item [Parámetros:] \hfill 
   \begin{description}
   \item[fpointer:] Puntero a fichero. 
   \end{description}
\item[Valores devueltos:] \hfill \\
   Devuelve un la posición del puntero al fichero como un número que representa el
   número de bytes desde el inicio del mismo. 
\end{description}
\end{framed}%
\lthtmlfigureZ
\lthtmlcheckvsize\clearpage}

\stepcounter{paragraph}
{\newpage\clearpage
\lthtmlfigureA{framed1506}%
\begin{framed}
\hfill \\$string$\  \textbf{fread} ( $mixed$\  \textbf{file}  )  
\begin{description}
\item [Parámetros:] \hfill 
   \begin{description}
   \item[file:] Fichero que será leído. Puede ser dado con ruta relativa al programa o absoluta, o como un puntero a fichero. Si el fichero no existe
   devolverá el valor nulo.
   \end{description}
\item[Valores devueltos:] \hfill \\
   Devuelve una cadena de caracteres que representa el contenido del fichero.
\end{description}
\end{framed}%
\lthtmlfigureZ
\lthtmlcheckvsize\clearpage}

\stepcounter{paragraph}
{\newpage\clearpage
\lthtmlfigureA{framed1517}%
\begin{framed}
\hfill \\$string$\  \textbf{fwrite} ( $mixed$\  \textbf{file}, $string$\  \textbf{content}  )  
\begin{description}
\item [Parámetros:] \hfill 
   \begin{description}
   \item[file:] Fichero que será leído. Puede ser dado con ruta relativa al programa o absoluta, o como un puntero a fichero. 
   Si el fichero no existe será creado.
   \item[content:] Cadena de caracteres que será el contenido del fichero.
   \end{description}
\item[Valores devueltos:] \hfill \\
   Devuelve la cadena de caracteres escrita en el fichero.
\end{description}
\end{framed}%
\lthtmlfigureZ
\lthtmlcheckvsize\clearpage}

\stepcounter{paragraph}
{\newpage\clearpage
\lthtmlfigureA{framed1529}%
\begin{framed}
\hfill \\$string$\  \textbf{fappend} ( $mixed$\  \textbf{file}, $string$\  \textbf{content}  )  
\begin{description}
\item [Parámetros:] \hfill 
   \begin{description}
   \item[file:] Fichero que será leído. Puede ser dado con ruta relativa al programa o absoluta, o como un puntero a fichero. 
   Si el fichero no existe será creado.
   \item[content:] Cadena de caracteres que será añadida al contenido del fichero.
   \end{description}
\item[Valores devueltos:] \hfill \\
   Devuelve la cadena de caracteres escrita en el fichero.
\end{description}
\end{framed}%
\lthtmlfigureZ
\lthtmlcheckvsize\clearpage}

\stepcounter{subsubsection}
\stepcounter{paragraph}
{\newpage\clearpage
\lthtmlfigureA{framed1542}%
\begin{framed}
\hfill \\$string$\  \textbf{date} ( $mixed$\  \textbf{format} [, $numeric$\  \textbf{timestamp} ] )  
\begin{description}
\item [Parámetros:] \hfill 
   \begin{description}
   \item[format:] Cadena de formato. Establece el formato de la fecha mediante una 
   serie de directivas.
   \begin{description}
         \item[\%d:] Día del mes con dos dígitos.
			\item[\%j:] Día del mes sin ceros iniciales.
			\item[\%l:] Día de la semana de forma alfabética completa.
			\item[\%D:] Día de la semana de forma alfabética y con tres letras.
			\item[\%w:] Día de la semana de forma numérica (0-domingo,6-sábado).
			\item[\%z:] Día del año de forma numérica.
			\item[\%F:] Mes de forma alfabética.
			\item[\%m:] Mes de forma numérica con dos dígitos.
			\item[\%n:] Mes de forma numérica sin ceros iniciales.
			\item[\%M:] Mes de forma alfabética con tres letras.
			\item[\%Y:] Año con cuatro dígitos.
			\item[\%y:] Año con dos dígitos.
			\item[\%a:] Periodo del día (am/pm) en minúsculas.
			\item[\%A:] Periodo del día (am/pm) en mayúsculas.
			\item[\%g:] Hora en formato 12h sin ceros iniciales.
			\item[\%G:] Hora en formato 24h sin ceros iniciales.
			\item[\%h:] Hora en formato 12h con dos dígitos.
			\item[\%H:] Hora en formato 24h con dos dígitos.
			\item[\%i:] Minutos con dos dígitos.
			\item[\%U:] Segundos desde la Época Unix (1 de Enero del 1970 00:00:00 GMT).
			\item[\%\%:] Carácter \%.
   \end{description}
   \item[timestamp:] Número entero que representa la marca de tiempo Unix a la que se desea dar formato. Si 
   no es dado se toma el momento actual.
   \end{description}
\item[Valores devueltos:] \hfill \\
   Devuelve la cadena de caracteres correspondiente a la marca de tiempo \textbf{timestamp} formateada según
   \textbf{format}.
\end{description}
\end{framed}%
\lthtmlfigureZ
\lthtmlcheckvsize\clearpage}

\stepcounter{paragraph}
{\newpage\clearpage
\lthtmlfigureA{framed1558}%
\begin{framed}
\hfill \\$numeric$\  \textbf{time} ( )  
\begin{description}
\item[Valores devueltos:] \hfill \\
   Devuelve el entero correspondiente a la marca de tiempo Unix actual.
\end{description}
\end{framed}%
\lthtmlfigureZ
\lthtmlcheckvsize\clearpage}

\stepcounter{subsection}
\stepcounter{subsubsection}
{\newpage\clearpage
\lthtmlinlinemathA{tex2html_wrap_inline4529}%
$ load$%
\lthtmlinlinemathZ
\lthtmlcheckvsize\clearpage}

\stepcounter{subsubsection}
{\newpage\clearpage
\lthtmlinlinemathA{tex2html_wrap_inline4534}%
$ dev$%
\lthtmlinlinemathZ
\lthtmlcheckvsize\clearpage}

{\newpage\clearpage
\lthtmlinlinemathA{tex2html_wrap_inline4536}%
$ run$%
\lthtmlinlinemathZ
\lthtmlcheckvsize\clearpage}

{\newpage\clearpage
\lthtmlinlinemathA{tex2html_wrap_inline4542}%
$ getParams$%
\lthtmlinlinemathZ
\lthtmlcheckvsize\clearpage}

\stepcounter{subsection}
\stepcounter{subsection}
\stepcounter{subsection}

\end{document}
